% English translation of chapter2.tex lines 519--720.
% Source: Methods of Algebra, Volume 2, commit
% 9a5803ff77dd3257484cb177f851a73770a59dd3 (CC BY 4.0).
% stable-unit-id: o014.aljabr2.chapter2.glimpse-of-lattice-theory

% segment-id: o014.li2.chapter2.glimpse-of-lattice-theory.h001
\section{A Glimpse of Lattice Theory}\label{sec:lattices}
\hypertarget{o014.aljabr2.chapter2.glimpse-of-lattice-theory}{ }

% segment-id: o014.li2.chapter2.glimpse-of-lattice-theory.p001
This section gives a brief introduction to a class of partially ordered
structures called lattices. The discussion will help clarify the structure of
abelian categories.

% segment-id: o014.li2.chapter2.glimpse-of-lattice-theory.p002
Throughout this section we consider various partially ordered sets $(P,\leq)$;
the corresponding category is denoted by $\mathcal{P}$.

% segment-id: o014.li2.chapter2.glimpse-of-lattice-theory.p003
For every subset $S$ of a partially ordered set $(P,\leq)$, there are notions
of upper and lower bounds, a supremum $\sup S$, and an infimum $\inf S$. A
subset $S\subset P$ has a supremum (or infimum) if and only if the coproduct
(or product) of all the objects in $S$ exists in the category $\mathcal{P}$;
in that case the supremum (or infimum) is precisely that coproduct (or
product).

% segment-id: o014.li2.chapter2.glimpse-of-lattice-theory.d001
\begin{definition}\label{def:lattice-order}
	\index{lattice}
	\index{lattice!bounded}
	\index{interval}
	\index[sym1]{veewedge@$\vee$, $\wedge$}
	If every two elements $a,b$ of a partially ordered set $(P,\leq)$ have a
	supremum $a\vee b$ and an infimum $a\wedge b$, then $(P,\leq)$ is called a
	\emph{lattice}. If, in addition, $(P,\leq)$ itself has an upper bound and a
	lower bound, they are both unique and are denoted by $1$ and $0$,
	respectively; in this case $P$ is called a \emph{bounded lattice}.

	If $a,b$ are elements of a partially ordered set $(P,\leq)$ and $a\leq b$,
	then $[a,b]:=\left\{x\in P:a\leq x\leq b\right\}$ is also a partially
	ordered set under $\leq$ and is called the \emph{interval} determined by
	$a,b$.
\end{definition}

% segment-id: o014.li2.chapter2.glimpse-of-lattice-theory.p004
If $P$ is a lattice, every interval in $P$ is closed under $\vee$ and
$\wedge$ and is a bounded lattice. If $P$ is a bounded lattice, then
$P=[0,1]$, and for every $x\in P$ one has
$x\wedge0=0$, $x\vee0=x=x\wedge1$, and $x\vee1=1$.

% segment-id: o014.li2.chapter2.glimpse-of-lattice-theory.p005
From the categorical point of view, $a\wedge b$, $a\vee b$, $1$, and $0$
correspond, respectively, to the product of two objects, the coproduct of two
objects, a terminal object (that is, an empty product), and an initial object
(that is, an empty coproduct) in the category; this can be checked directly.

% segment-id: o014.li2.chapter2.glimpse-of-lattice-theory.p006
As an illustration of the operations $\wedge$ and $\vee$, prove the following
fact. If $a^\flat,a,b\in P$ satisfy $a^\flat\leq a$, then
% segment-id: o014.li2.chapter2.glimpse-of-lattice-theory.q001
\begin{equation}\label{eqn:modular-inequality}
	a^\flat \vee (a \wedge b) \leq (a^\flat \vee b) \wedge a.
\end{equation}
By the definition of a supremum, it first suffices to prove
$a^\flat\leq(a^\flat\vee b)\wedge a$ and
$(a\wedge b)\leq(a^\flat\vee b)\wedge a$. By the definition of an infimum,
it then suffices to prove, respectively, the inequalities
\[ a^\flat \leq a^\flat \vee b, \quad a^\flat \leq a, \quad a \wedge b \leq a^\flat \vee b, \quad a \wedge b \leq a; \]
the third follows from $a\wedge b\leq b\leq a^\flat\vee b$, and the others
are clear.

% segment-id: o014.li2.chapter2.glimpse-of-lattice-theory.d002
\begin{definition}\label{def:modular-lattice}
	\index{lattice!modular}
	\index{complement}
	Let $P$ be a lattice.
	% segment-id: o014.li2.chapter2.glimpse-of-lattice-theory.l001
	\begin{itemize}
		% segment-id: o014.li2.chapter2.glimpse-of-lattice-theory.i001
		\item $P$ is called a \emph{modular lattice} if the following property
		holds: whenever $a^\flat,a,b\in P$ satisfy $a^\flat\leq a$, then
		% segment-id: o014.li2.chapter2.glimpse-of-lattice-theory.q002
		\[ a^\flat \vee (a \wedge b) = (a^\flat \vee b) \wedge a. \]
		% segment-id: o014.li2.chapter2.glimpse-of-lattice-theory.i002
		\item Suppose that $P$ is a bounded lattice. If $x,c\in P$ satisfy
		$x\vee c=1$ and $x\wedge c=0$, then $c$ is called a \emph{complement} of
		$x$ in $P$.
	\end{itemize}
\end{definition}

% segment-id: o014.li2.chapter2.glimpse-of-lattice-theory.p007
We now consider some basic examples of lattices.
% segment-id: o014.li2.chapter2.glimpse-of-lattice-theory.l002
\begin{itemize}
	% segment-id: o014.li2.chapter2.glimpse-of-lattice-theory.i003
	\item The set of positive integers $\Z_{\geq1}$ is a lattice under
	divisibility ($x\mid y\iff x\leq y$):
	$x\vee y=\mathrm{lcm}(x,y)$ and $x\wedge y=\mathrm{gcd}(x,y)$. The reader
	can check that this lattice is modular, has lower bound $1$, but has no
	upper bound.

	% segment-id: o014.li2.chapter2.glimpse-of-lattice-theory.i004
	\item Let $R$ be a ring and $M$ a left $R$-module. Its set of submodules
	$\mathrm{Sub}_M$ is a bounded modular lattice under $\subset$:
	$x\vee y=x+y$ and $x\wedge y=x\cap y$; its upper bound is $M$ and its lower
	bound is $\{0\}$. This is the origin of the term ``modular lattice.''

	% segment-id: o014.li2.chapter2.glimpse-of-lattice-theory.i005
	\item All closed subspaces of a Hilbert space $H$ form a bounded lattice
	under $\subset$: $x\vee y:=\overline{x+y}$ and $x\wedge y=x\cap y$. One can
	prove that this lattice is not modular when $H$ is infinite-dimensional
	(Birkhoff--von Neumann).
	% Possible reference: Proposition 4.4 in https://link.springer.com/chapter/10.1007/978-94-015-9026-6_4
\end{itemize}

% segment-id: o014.li2.chapter2.glimpse-of-lattice-theory.r001
\begin{remark}\label{rem:lattice-duality}
	The definition of a lattice also contains a duality. Given the partial order
	$\leq$ on a set $P$, we may consider the opposite partial order
	$\leq^{\opp}$: $x\leq^{\opp}y\iff x\geq y$; this is equivalent to replacing
	$\mathcal{P}$ by $\mathcal{P}^{\opp}$. If $(P,\leq)$ is a lattice (or a
	bounded lattice), then so is $(P,\leq^{\opp})$. Moreover, $\wedge,\vee,0,1$
	in $(P,\leq)$ correspond, respectively, to $\vee,\wedge,1,0$ in
	$(P,\leq^{\opp})$.
\end{remark}

% segment-id: o014.li2.chapter2.glimpse-of-lattice-theory.le001
\begin{lemma}\label{prop:modularity-criterion}
	Let $P$ be a lattice. Then $P$ is modular if and only if the following
	property holds for every interval $I$ in $P$: if $c^\flat,c$ are both
	complements of $x\in I$ in $I$ and $c^\flat\leq c$, then $c^\flat=c$.
\end{lemma}
% segment-id: o014.li2.chapter2.glimpse-of-lattice-theory.pf001
\begin{proof}
	If $P$ is a modular lattice, then every interval $I$ is also a modular
	lattice. Without loss of generality, assume $I=P$; for the $x,c^\flat,c$ in
	the statement, we have
	% segment-id: o014.li2.chapter2.glimpse-of-lattice-theory.q003
	\[ c = c \wedge 1 = c \wedge (x \vee c^\flat) = c^\flat \vee (x \wedge c) = c^\flat \vee 0 = c^\flat . \]

	Conversely, assume that the condition on complements holds for every
	interval. Let $a^\flat,a,b\in P$ satisfy $a^\flat\leq a$. Set
	% segment-id: o014.li2.chapter2.glimpse-of-lattice-theory.q004
	\[ b \wedge a \leq c_1 := a^\flat \vee (b \wedge a) \underset{\because \eqref{eqn:modular-inequality}}{\leq} a \wedge (b \vee a^\flat) =: c_2 \leq b \vee a^\flat. \]
	We prove that $c_1,c_2$ are both complements of $b$ in
	$\left[b\wedge a,\;b\vee a^\flat\right]$, so that $c_1=c_2$. First, the
	definition immediately gives $c_1\leq a$; therefore
	% segment-id: o014.li2.chapter2.glimpse-of-lattice-theory.q005
	\[ b \wedge a \geq b \wedge c_1 = (a^\flat \vee (b \wedge a)) \wedge b \underset{\because \eqref{eqn:modular-inequality}}{\geq} (a^\flat \wedge b) \vee (b \wedge a) = b \wedge a , \]
	so $b\wedge c_1=b\wedge a$. On the other hand,
	$b\vee c_1=a^\flat\vee(b\wedge a)\vee b=a^\flat\vee b$. Thus $c_1$ is
	indeed a complement of $b$ in
	$\left[b\wedge a,\;b\vee a^\flat\right]$.

	Lattice duality (Remark~\ref{rem:lattice-duality}) reverses the order
	relation between $a^\flat,a$, interchanges $\wedge$ and $\vee$, and
	interchanges the roles of $c_1$ and $c_2$. Thus the preceding argument,
	carried out in $(P,\leq^{\opp})$, shows that $c_2$ is a complement of $b$
	in $\left[b\wedge a,\;b\vee a^\flat\right]$.
\end{proof}

% segment-id: o014.li2.chapter2.glimpse-of-lattice-theory.pr001
\begin{proposition}\label{prop:lattice-diamond-isom}
	\index{standard isomorphism}
	Let $a,b$ be elements of a modular lattice $(P,\leq)$. Then there is an
	isomorphism of partially ordered sets
	% segment-id: o014.li2.chapter2.glimpse-of-lattice-theory.g001
	\[\begin{tikzcd}[row sep=tiny]
		{[a \wedge b, a]} \arrow[leftrightarrow, r, "1:1"] & {[b, a \vee b]} \\
		x \arrow[mapsto, r] & x \vee b \\
		y \wedge a & y \arrow[mapsto, l]
	\end{tikzcd}\]
	called the \emph{standard isomorphism}.
\end{proposition}
% segment-id: o014.li2.chapter2.glimpse-of-lattice-theory.pf002
\begin{proof}
	The two maps are clearly well defined and order-preserving. For
	$x\in[a\wedge b,a]$, the definition of a modular lattice gives
	$(x\vee b)\wedge a=(a\wedge b)\vee x$, whose right-hand side is $x$.
	Dually, $y\in[b,a\vee b]$ implies $(y\wedge a)\vee b=y$
	(Remark~\ref{rem:lattice-duality}); thus the two maps are mutually inverse.
\end{proof}

% segment-id: o014.li2.chapter2.glimpse-of-lattice-theory.p008
The Zassenhaus lemma of elementary algebra (see \cite[Lemma 4.6.4]{Li1},
especially its module-theoretic version) and its consequences can be distilled
using lattice theory.

% segment-id: o014.li2.chapter2.glimpse-of-lattice-theory.th001
\begin{theorem}[Zassenhaus Lemma for Modular Lattices]\label{prop:lattice-Zassenhaus}
	\index{Zassenhaus Lemma}
	Let $(P,\leq)$ be a modular lattice. For given elements
	$u^\flat\leq u$ and $v^\flat\leq v$, there is a diagram
	% segment-id: o014.li2.chapter2.glimpse-of-lattice-theory.g002
	\[ \begin{tikzcd}[column sep=tiny, row sep=small]
		u^\flat \vee (u \wedge v) \arrow[dash, dd] \arrow[dash, rd] & & (u \wedge v) \vee v^\flat \arrow[dash, dd] \arrow[dash, ld] \\
		& u \wedge v \arrow[dash, dd] & \\
		u^\flat \vee (u \wedge v^\flat) \arrow[dash, rd] & & (u^\flat \wedge v) \vee v^\flat \arrow[dash, ld] \\
		& (u^\flat \wedge v) \vee (u \wedge v^\flat)
	\end{tikzcd} \]
	meaning that its elements satisfy the following patterns:
	% segment-id: o014.li2.chapter2.glimpse-of-lattice-theory.q006
	\begin{equation}\label{eqn:Zassenhaus-cases}
	% segment-id: o014.li2.chapter2.glimpse-of-lattice-theory.g003
	\begin{tikzcd}
		x \arrow[dash, d, "\geq" description, sloped] \\ y
	\end{tikzcd} \qquad
	% segment-id: o014.li2.chapter2.glimpse-of-lattice-theory.g004
	\begin{tikzcd}[column sep=tiny]
		x \arrow[dash, rd] & & y \arrow[dash, ld] \\
		& x \wedge y &
	\end{tikzcd} \quad
	% segment-id: o014.li2.chapter2.glimpse-of-lattice-theory.g005
	\begin{tikzcd}[column sep=tiny]
		{} & x \vee y \arrow[dash, ld] \arrow[dash, rd] & \\
		x & & y
	\end{tikzcd} .
	\end{equation}
	Moreover, there are standard isomorphisms between the intervals
	% segment-id: o014.li2.chapter2.glimpse-of-lattice-theory.q007
	\begin{multline*}
		\left[ u^\flat \vee (u \wedge v^\flat), \; u^\flat \vee (u \wedge v) \right] \leftiso \left[ (u^\flat \wedge v) \vee (u \wedge v^\flat), \; u \wedge v \right] \\
		\rightiso \left[ (u^\flat \wedge v) \vee v^\flat, \; (u \wedge v) \vee v^\flat \right],
	\end{multline*}
	realized by the maps
	$x\vee u^\flat\vee(u\wedge v^\flat)\mapsfrom x\mapsto
	x\vee(u^\flat\wedge v)\vee v^\flat$ in
	Proposition~\ref{prop:lattice-diamond-isom}.
\end{theorem}
% segment-id: o014.li2.chapter2.glimpse-of-lattice-theory.pf003
\begin{proof}
	The isomorphisms in the last part are obtained by applying
	Proposition~\ref{prop:lattice-diamond-isom} to the two parallelograms in the
	diagram; what remains to be checked are the relations in
	\eqref{eqn:Zassenhaus-cases}. Notice that the diagram is left--right
	symmetric under $u\leftrightarrow v$ and
	$u^\flat\leftrightarrow v^\flat$; it therefore suffices to check its left
	half. The partial-order relations among its elements are immediate. For the
	$\wedge$ case, the definition of a modular lattice gives
	% segment-id: o014.li2.chapter2.glimpse-of-lattice-theory.q008
	\begin{equation*}
		\left( u^\flat \vee (u \wedge v^\flat)\right) \wedge (u \wedge v) = \left( u^\flat \wedge (u \wedge v) \right) \vee (u \wedge v^\flat) = (u^\flat \wedge v) \vee (u \wedge v^\flat).
	\end{equation*}
	For the $\vee$ case, the inequality $u\wedge v^\flat\leq u\wedge v$
	immediately gives
	% segment-id: o014.li2.chapter2.glimpse-of-lattice-theory.q009
	\begin{equation*}
		u^\flat \vee (u \wedge v^\flat) \vee (u \wedge v) = u^\flat \vee (u \wedge v).
	\end{equation*}
	This proves \eqref{eqn:Zassenhaus-cases}.
\end{proof}

% segment-id: o014.li2.chapter2.glimpse-of-lattice-theory.p009
Next consider a finite descending chain in a lattice, written
$x_0\geq x_1\geq\cdots\geq x_r$ ($r\in\Z_{\geq0}$). A descending chain
obtained by inserting finitely many intermediate elements is called a
\emph{refinement} of the original descending chain; if the inserted elements
include some $y\notin\{x_0,\ldots,x_r\}$, it is called a proper refinement.

% segment-id: o014.li2.chapter2.glimpse-of-lattice-theory.d003
\begin{definition}\label{def:Schreier-equiv}
	Two descending chains of the same length in a modular lattice,
	$x_0\geq\cdots\geq x_r$ and $x'_0\geq\cdots\geq x'_r$, are called
	\emph{equivalent} if the following condition holds: there is a bijection
	$\sigma$ from $\{0,\ldots,r-1\}$ to itself (that is, a
	permutation)\footnote{Translator's note (O014-C021): the source has
	$\{0,\ldots,r\}$, but the successive factors being paired are indexed by
	$0\leq i<r$; the index $i=r$ would leave $x_{i+1}$ undefined.} such that,
	for every $0\leq i<r$, there is an isomorphism of partially ordered sets
	% segment-id: o014.li2.chapter2.glimpse-of-lattice-theory.q010
	\[ \tau_i: \left[ x_{i+1}, x_i \right] \simeq \left[ x'_{\sigma(i) + 1}, x'_{\sigma(i)} \right], \]
	and $\tau_i$ decomposes into a finite composite of standard isomorphisms
	(Proposition~\ref{prop:lattice-diamond-isom}) or their inverses.
\end{definition}

% segment-id: o014.li2.chapter2.glimpse-of-lattice-theory.th002
\begin{theorem}[Schreier Refinement Theorem for Modular Lattices]\label{prop:Schreier-refinement}
	\index{Schreier Refinement Theorem}
	\mbox{}\par
	Consider descending chains
	% segment-id: o014.li2.chapter2.glimpse-of-lattice-theory.q011
	\[ x_0 \geq \cdots \geq x_r, \quad y_0 \geq \cdots \geq y_s, \]
	in a modular lattice $(P,\leq)$ such that $x_0=y_0$ and $x_r=y_s$, with
	$r,s\in\Z_{\geq0}$. Then each chain has a refinement, and the two
	refinements are equivalent.
\end{theorem}
% segment-id: o014.li2.chapter2.glimpse-of-lattice-theory.pf004
\begin{proof}
	The argument is the same as in the case of groups
	\cite[Theorem 4.6.6]{Li1}; both are based on
	Theorem~\ref{prop:lattice-Zassenhaus}. We give the outline. Define
	% segment-id: o014.li2.chapter2.glimpse-of-lattice-theory.q012
	\[ x_{i, j} := x_{i+1} \vee (x_i \wedge y_j), \quad y_{j, i} := (x_i \wedge y_j) \vee y_{j+1}, \]
	with $0\leq i<r$ and $0\leq j\leq s$ for $x_{i,j}$, and with
	$0\leq i\leq r$ and $0\leq j<s$ for $y_{j,i}$. Then
	$x_{i,j+1}\leq x_{i,j}$, $x_{i,0}=x_i$, and $x_{i,s}=x_{i+1}$, so
	$\left(x_{i,j}\right)_{i,j}$, in that order, refines
	$x_0\geq\cdots\geq x_r$. Similarly,
	$\left(y_{j,i}\right)_{j,i}$ refines $y_0\geq\cdots\geq y_s$. Taking
	$u^\flat=x_{i+1}$, $u=x_i$, $v^\flat=y_{j+1}$, and $v=y_j$ in
	Theorem~\ref{prop:lattice-Zassenhaus}, we obtain
	% segment-id: o014.li2.chapter2.glimpse-of-lattice-theory.q013
	\[ \left[ x_{i, j+1}, x_{i,j} \right] \simeq \left[ y_{j, i+1}, y_{j, i} \right], \]
	and this isomorphism decomposes into standard isomorphisms and their
	inverses. The assertion follows.
\end{proof}

% segment-id: o014.li2.chapter2.glimpse-of-lattice-theory.p010
From now on, we consider only strictly ascending or strictly descending
chains.

% segment-id: o014.li2.chapter2.glimpse-of-lattice-theory.d004
\begin{definition}\index{composition series}
	Fix a partially ordered set $P$ and elements $a\leq b$ in
	it.\footnote{Translator's note (O014-C022): the source has $a<b$, but the
	next sentence explicitly includes the case $a=b$ and $r=0$; the target uses
	the weaker hypothesis consistent with that case.} If a descending chain
	$b=x_0>\cdots>x_r=a$ in $P$ has no proper refinement, it is called a
	\emph{composition series} of $[a,b]$.
\end{definition}

% segment-id: o014.li2.chapter2.glimpse-of-lattice-theory.p011
The length of this composition series is defined to be
$r\in\Z_{\geq0}$; notice that $r=0$ if and only if $a=b$.
\index{length}

% segment-id: o014.li2.chapter2.glimpse-of-lattice-theory.th003
\begin{theorem}[Jordan--Hölder Theorem for Modular Lattices]\label{prop:lattice-JH}
	\index{Jordan-Holder@Jordan--Hölder Theorem}
	Let $a<b$ be elements of a modular lattice $(P,\leq)$. Then any two
	composition series of $[a,b]$ have the same length and are equivalent in
	the sense of Definition~\ref{def:Schreier-equiv}.
\end{theorem}
% segment-id: o014.li2.chapter2.glimpse-of-lattice-theory.pf005
\begin{proof}
	This follows immediately from
	Theorem~\ref{prop:Schreier-refinement}.
\end{proof}

% segment-id: o014.li2.chapter2.glimpse-of-lattice-theory.p012
The next important question is to determine which intervals in a partially
ordered set have composition series. As in the familiar theory of modules,
their existence can be guaranteed by ascending- and descending-chain
conditions.

% segment-id: o014.li2.chapter2.glimpse-of-lattice-theory.d005
\begin{definition}\label{def:lattice-Noether-Artin}
	\index{partially ordered set!finite length}
	\index{partially ordered set!Artinian, Noetherian}
	Let $(P,\leq)$ be a partially ordered set. If there is no infinite
	ascending chain $x_1<x_2<x_3<\cdots$ in $P$, then $P$ is called
	\emph{Noetherian}; if there is no infinite descending chain
	$x_1>x_2>x_3>\cdots$, then $P$ is called \emph{Artinian}. A partially
	ordered set that is both Artinian and Noetherian is said to have
	\emph{finite length}.
\end{definition}

% segment-id: o014.li2.chapter2.glimpse-of-lattice-theory.p013
If $(P,\leq)$ is Noetherian (or Artinian, or of finite length), then every
subset of it has the same property. It is easy to prove that the Noetherian
(or Artinian) condition is equivalent to the assertion that every nonempty
subset of $P$ has a maximal (or minimal) element with respect to $\leq$. In
this chapter these notions will be applied chiefly to partially ordered sets
of subobjects; see Definition~\ref{def:sub-quot-obj}.

% segment-id: o014.li2.chapter2.glimpse-of-lattice-theory.d006
\begin{definition}\label{def:finite-length-object}
	\index{object!finite length}
	\index{object!Noetherian}
	\index{object!Artinian}
	Fix a category $\mathcal{C}$. An object $X$ of $\mathcal{C}$ is called
	Noetherian (or Artinian, or of finite length) if its partially ordered set
	of subobjects $(\mathrm{Sub}_X,\subset)$ is Noetherian (or Artinian, or of
	finite length).
\end{definition}

% segment-id: o014.li2.chapter2.glimpse-of-lattice-theory.p014
We now return to general partially ordered sets.

% segment-id: o014.li2.chapter2.glimpse-of-lattice-theory.le002
\begin{lemma}\label{prop:lattice-composition-series}
	Let $a\leq b$ be elements of a partially ordered set $(P,\leq)$.
	% segment-id: o014.li2.chapter2.glimpse-of-lattice-theory.l003
	\begin{enumerate}[(i)]
		% segment-id: o014.li2.chapter2.glimpse-of-lattice-theory.i006
		\item If $[a,b]$ has finite length, then every chain in it can be refined
		to a composition series; in particular, $[a,b]$ has a composition series.
		% segment-id: o014.li2.chapter2.glimpse-of-lattice-theory.i007
		\item If $[a,b]$ is a modular lattice and has a composition series, then
		$[a,b]$ has finite length.
	\end{enumerate}
\end{lemma}
% segment-id: o014.li2.chapter2.glimpse-of-lattice-theory.pf006
\begin{proof}
	Without loss of generality, assume $a<b$. For (i), given a chain
	$y_0>\cdots>y_r$, it suffices to show that every $[y_{i+1},y_i]$ has a
	composition series.\footnote{Translator's note (O014-C023): the source has
	$[y_i,y_{i+1}]$ and then sets $y_i=a$, $y_{i+1}=b$. Since
	$y_i>y_{i+1}$, both orders are reversed; the target uses the interval
	$[y_{i+1},y_i]$ and the well-typed assignment.} Without loss of generality,
	assume $y_{i+1}=a$ and $y_i=b$. The Artinian condition guarantees a minimal
	element $x^0\in[a,b]$ with $x^0>a$. Similarly, if $x^0\neq b$, there is a
	minimal element $x^1\in[a,b]$ with $x^1>x^0$, and so on. By the Noetherian
	condition, this process must stop after finitely many steps, producing a
	chain $b>\cdots>x^0>a$ with no proper refinement.

	For (ii), choose a composition series $b=x_0>\cdots>x_r=a$.
	Theorem~\ref{prop:Schreier-refinement} shows that every finite chain
	$\cdots>y_i>y_{i+1}>\cdots$ has a refinement equivalent to the composition
	series above; hence the length of that chain cannot exceed the length of the
	composition series. Thus $[a,b]$ is a partially ordered set of finite
	length.
\end{proof}

% segment-id: o014.li2.chapter2.glimpse-of-lattice-theory.d007
\begin{definition}\label{def:lattice-length}
	Let $(P,\leq)$ be a bounded modular lattice of finite length. Choose any
	composition series $1=x_0>\cdots>x_r=0$ of $P$. The number
	$r\in\Z_{\geq0}$ is called the \emph{length} of $(P,\leq)$.
\end{definition}

% segment-id: o014.li2.chapter2.glimpse-of-lattice-theory.p015
Theorem~\ref{prop:lattice-JH} shows that this length is well defined and does
not depend on the chosen composition series; the length of $(P,\leq)$ is $0$
if and only if $P$ is a one-point set.
